\section{積分方程式}\label{節:積分方程式}
未知の関数が積分の中にも現れる等式を,\textbf{積分方程式}という.積分の中にある未知の関数を,どうすれば取り出せるだろうか.上端が変数なら,前節までの結果から,両辺を微分すれば積分を外せる.微分方程式に帰着させるのである.

ただし,微分すると定数の情報が失われる.たとえば $f(x)=2+\int_0^x f(t)\dd t$ を微分すると,定数 $2$ は消えてしまう.そこで,微分する前に積分区間の両端が一致する値を代入して $f$ の値を取り出しておき,最後に求めた関数を元の式へ代入して確かめる.

\begin{bookpoint}
\textbf{POINT}\par
積分区間の両端が一致する $x$ を代入して初期値を得る.両辺を微分して微分方程式にし,解いた関数を元の式へ代入して確かめる.
\end{bookpoint}

\begin{example}{積分方程式}
連続関数 $f$ が,すべての実数 $x$ について
\[
 f(x)=2+\int_0^x f(t)\dd t
\]
を満たすとき,\,$f$ を求めよ.
\end{example}
\paragraphbf{方針}
$x=0$ を代入して $f(0)$ を求め,両辺を微分して $f'=f$ に帰着させる.\,$f'=f$ は\sectionref{節:微分方程式}と同じく $\e^x$ で割って解く.

\paragraphbf{解答}
$f(x)=2\e^x$

\paragraphbf{解法}
$x=0$ を代入すると,\,$\int_0^0f(t)\dd t=0$ より
\[
 f(0)=2
\]
である.\,$f$ が連続なので,右辺の積分は基本定理により微分可能である.したがって左辺の $f$ も微分可能となり,両辺を微分して
\[
 f'(x)=f(x)
\]
を得る.

この等式だけで直ちに $f=\e^x$ としてはいけない.定数倍も同じ関係を満たすからである.\,$f$ が $\e^x$ の定数倍なら,\,$\e^x$ で割った $\e^{-x}f(x)$ は定数になるはずである.そこで,この商を微分すると,
\begin{align*}
 \frac{\dd}{\dd x}\{\e^{-x}f(x)\}
 &=\e^{-x}\{f'(x)-f(x)\}\\
 &=0
\end{align*}
となる.よって $\e^{-x}f(x)$ は定数であり,\,$x=0$ での値からその定数は $2$ である.したがって候補は $f(x)=2\e^x$ となる.

最後に元の式へ代入して確かめる.
\begin{align*}
 2+\int_0^x2\e^t\dd t
 &=2+\left[2\e^t\right]_0^x\\
 &=2+2\e^x-2\\
 &=2\e^x
\end{align*}
確かに成り立つので,求める関数は $f(x)=2\e^x$ である.微分して得た条件は必要条件にすぎないため,この確認によって候補が実際に解であることを示した.

ここまでで,定積分の意味と,その計算を微分に結び付ける根拠がそろった.次章からは,原始関数を見つけるための式の変形を詳しく扱う.どの公式を使う場合も,求めた原始関数を微分して被積分関数に戻ることを確かめられる.
